% Part: history
% Chapter: biographies
% Section: julia-robinson

\documentclass[../../../include/open-logic-section]{subfiles}

\begin{document}

\olfileid{his}{bio}{rob}

\olsection{জুলিয়া রবিনসন}

\olphoto{robinson-julia}{জুলিয়া রবিনসন}

জুলিয়া বোম্যান রবিনসন ছিলেন একজন মার্কিন গণিতবিদ।
সিদ্ধান্ত-সমস্যা নিয়ে তাঁর কাজের জন্য তিনি প্রধানত
পরিচিত; বিশেষত হিলবার্টের দশম সমস্যার সমাধানে
তাঁর অবদান বিখ্যাত। ১৯১৯ সালের ৮ ডিসেম্বর
মিসৌরির সেন্ট লুইসে তাঁর জন্ম। রবিনসনের স্মৃতিতে
ছোটবেলা থেকেই সংখ্যার প্রতি তাঁর কৌতূহল ছিল
\citep[4]{Reid1986}। নয় বছর বয়সে তাঁর
স্কারলেট জ্বর হয়; তার পরে কয়েক বার বাতজ্বর
ফিরে আসে। এতে তাঁকে অনেকটা সময় বিছানায়
কাটাতে হয় এবং পড়াশোনায় তিনি পিছিয়ে পড়েন।
গৃহশিক্ষকদের সাহায্যে পরে ঘাটতি পূরণ করতে
পারলেও অসুস্থতার শারীরিক প্রভাব তাঁর জীবনে
দীর্ঘস্থায়ী হয়।

শৈশবের এই সমস্যাগুলি সত্ত্বেও গণিত ও বিজ্ঞানে
কয়েকটি পুরস্কার নিয়ে রবিনসন উচ্চবিদ্যালয়ের
পড়া শেষ করেন। সান ডিয়েগো স্টেট কলেজে
বিশ্ববিদ্যালয়ের পড়া শুরু করে শেষ বর্ষে তিনি
ক্যালিফোর্নিয়া বিশ্ববিদ্যালয়ের বার্কলি শাখায়
যান। সেখানে গণিতবিদ রাফায়েল রবিনসনের
প্রভাবে আসেন। দুজনের বন্ধুত্ব হয় এবং ১৯৪১
সালে তাঁরা বিয়ে করেন। একজন অধ্যাপকের স্ত্রী
হওয়ায় বার্কলির গণিত বিভাগে রবিনসনের
পড়ানো নিষিদ্ধ ছিল। গণিতের ক্লাসে অনুমতি
নিয়ে বসা চালিয়ে গেলেও তিনি বিশ্ববিদ্যালয়
ছেড়ে সংসার শুরু করার আশা করেছিলেন।
কিন্তু বিয়ের কিছুদিন পরেই তাঁর নিউমোনিয়া
হয়। চিকিৎসকেরা জানান, শৈশবের বাতজ্বরের
কারণে তাঁর হৃদ্‌যন্ত্রে যথেষ্ট ক্ষতচিহ্ন জমেছে।
সেই ক্ষতির মাত্রা দেখে চিকিৎসক বলেন যে
তিনি চল্লিশ বছরের বেশি বাঁচবেন না; সন্তান
নিতেও তাঁকে বারণ করা হয় \citep[13]{Reid1986}।

অনেক দিন বিষণ্ণ থাকার পর রবিনসন আবার
গণিত নিয়ে পড়াশোনা করার সিদ্ধান্ত নেন।
বার্কলিতে ফিরে আলফ্রেড তার্স্কির তত্ত্বাবধানে
১৯৪৮ সালে পিএইচডি সম্পূর্ণ করেন। বাস্তব
সংখ্যার প্রথম-ক্রমের তত্ত্ব সিদ্ধান্তযোগ্য, এটি
তার্স্কি দেখিয়েছিলেন; আর গ্যোডেলের কাজ
থেকে পাওয়া গিয়েছিল যে স্বাভাবিক সংখ্যার
প্রথম-ক্রমের তত্ত্ব সিদ্ধান্তযোগ্য নয়। মূলদ
সংখ্যার প্রথম-ক্রমের তত্ত্ব সিদ্ধান্তযোগ্য কি না,
তা ছিল একটি বড় অমীমাংসিত প্রশ্ন। নিজের
অভিসন্দর্ভে \citeyearpar{Robinson1949}
রবিনসন প্রমাণ করেন যে সেটি সিদ্ধান্তযোগ্য নয়।

সিদ্ধান্ত-সমস্যায় আগ্রহী রবিনসন এরপর
হিলবার্টের দশম সমস্যার সমাধান খুঁজতে শুরু
করেন। ১৯০০ সালে ডাভিড হিলবার্টের উত্থাপিত
২৩টি বিখ্যাত গাণিতিক সমস্যার তালিকায়
এটি ছিল দশম। প্রশ্নটি হলো, এমন কোনো
অ্যালগরিদম আছে কি, যা সীমিত সময়ে
বলতে পারবে যে পূর্ণসংখ্যা সহগবিশিষ্ট
কোনো বহুপদী সমীকরণের—যেমন
$3x^2 - 2y +3 = 0$—পূর্ণসংখ্যায় সমাধান
আছে কি না। এই ধরনের প্রশ্নকে
\emph{ডায়োফ্যান্টাইন সমস্যা} বলে। প্রথমদিকের
কিছু সাফল্যের পরে রবিনসন মার্টিন ডেভিস
ও হিলারি পাটনামের সঙ্গে কাজ শুরু করেন;
তাঁরাও সমস্যাটি নিয়ে কাজ করছিলেন।
তাঁরা প্রমাণ করেন যে সূচকীয় ডায়োফ্যান্টাইন
সমস্যা—যেখানে অজ্ঞাত রাশি সূচকেও থাকতে
পারে—সিদ্ধান্তযোগ্য নয়। আরও দেখান, একটি
নির্দিষ্ট অনুমান (পরে “J.R.” নামে পরিচিত)
থেকে হিলবার্টের দশম সমস্যার অসিদ্ধান্তযোগ্যতা
পাওয়া যায় \citep{DavisPutnamRobinson1961}।
১৯৬০-এর দশকজুড়ে রবিনসন সমস্যাটি নিয়ে
কাজ চালিয়ে যান। ১৯৭০ সালে তরুণ রুশ
গণিতবিদ ইউরি মাতিয়াসেভিচ অবশেষে J.R.
অনুমানটি প্রমাণ করেন। সম্মিলিত ফলটি এখন
মাতিয়াসেভিচ--রবিনসন--ডেভিস--পাটনাম
উপপাদ্য, সংক্ষেপে MRDP উপপাদ্য নামে
পরিচিত। মাতিয়াসেভিচ ও রবিনসনের বন্ধুত্ব
হয় এবং তাঁরা একসঙ্গে কয়েকটি প্রবন্ধ লেখেন।
মাতিয়াসেভিচকে একটি চিঠিতে রবিনসন একবার
লিখেছিলেন, হাজার হাজার মাইল দূরে থেকেও
একসঙ্গে কাজ করে তাঁরা দুজনের আলাদা
প্রচেষ্টার চেয়ে বেশি এগোতে পারছেন বলে
তিনি খুবই খুশি \citep[45]{Matijasevich1992}।

রবিনসন ছিলেন আমেরিকান ম্যাথমেটিক্যাল
সোসাইটির প্রথম নারী সভাপতি এবং
ন্যাশনাল অ্যাকাডেমি অব সায়েন্সেসের
সদস্য নির্বাচিত প্রথম নারী। লিউকেমিয়া
ধরা পড়ার পর ১৯৮৫ সালের ৩০ জুলাই
৬৫ বছর বয়সে তাঁর মৃত্যু হয়।

\begin{reading}
রবিনসনের গাণিতিক প্রবন্ধগুলি তাঁর
\textit{Collected Works} (সংকলিত রচনাবলি)
\citep{Robinson1996}-তে পাওয়া যায়।
এতে ন্যাশনাল অ্যাকাডেমি অব সায়েন্সেসের
জন্য লেখা তাঁর জীবনীমূলক স্মারক রচনাটির
পুনর্মুদ্রণও আছে \citep{Feferman1994}।
রবিনসনের দিদি কনস্ট্যান্স রিড সাক্ষাৎকারের
ভিত্তিতে “Autobiography of Julia” প্রকাশ
করেন \citep{Reid1986}; একটি পূর্ণাঙ্গ
স্মৃতিকথাও লেখেন \citep{Reid1996}।
রবিনসন ও হিলবার্টের দশম সমস্যা নিয়ে
একটি স্বল্পদৈর্ঘ্যের প্রামাণ্যচিত্র পরিচালনা
করেন জর্জ চিচেরি \citep{Csicsery2016}।
মাতিয়াসেভিচের সঙ্গে রবিনসনের যৌথ কাজ
এবং মাতিয়াসেভিচের কাজে তাঁর প্রভাব
নিয়ে সংক্ষিপ্ত স্মৃতিকথার জন্য
\citep{Matijasevich1992} দেখো।
\end{reading}

\end{document}
